% !PRINTED 49
\weberpage{49}

Wenn umgekehrt einer der Charaktere $\chi = \xi$ der Gruppe $S$
für alle Elemente der Gruppe $T$ denselben Werth hat, so kann
dies nur der Werth 1 sein, da in $T$ sicher das Einheitselement
von $S$ vorkommt (\S.\,11, 2.). Wenn dann $A_i$ irgend ein Element
aus $T_i$ ist, und $A$ die ganze Gruppe $T$ durchläuft, so
durchläuft $AA_i$ die Nebengruppe $T_i$. Es ist dann aber
$\xi(A) = 1$ und folglich
\weq{5}{\xi(AA_i) = \xi(A_i),}
d.\,h. $\xi(A)$ hat für alle Elemente einer Nebengruppe $T_i$ einen
und denselben Werth, und $\xi(A)$ kann also als Function von $T_i$
aufgefasst und mit $\xi(T_i)$ bezeichnet werden.

Da nun, wenn $A_i$ in $T_i$ und $A_k$ in $T_k$ vorkommt, $A_iA_k$
in $T_iT_k$ enthalten ist, so folgt
\weq{6}{\xi(T_i)\xi(T_k) = \xi(T_iT_k),}
d.\,h. $\xi(T_i)$ ist unter den Charakteren der Gruppe $S/T$
enthalten. Es giebt also genau $j$ solche Charaktere $\xi(A)$, die
dadurch definirt sind, dass sie für jedes Element in $T$ den Werth
1 haben. Diese $j$ Charaktere $\xi$ bilden nach der Composition
der Charaktere eine Gruppe; denn ist $\xi_1(A) = 1$,
$\xi_2(A) = 1$, so ist auch $\xi_1\xi_2(A) = 1$ [\S.\,11, (18)].
Da die Functionen $\xi$ nach (4) auch als die Charaktere der
Gruppe $S/T$ aufgefasst werden können, so ist nach \S.\,11, 5. die
Gruppe der $\xi$ mit der Gruppe $S/T$ isomorph.

Damit ist also der folgende Satz bewiesen:

7.\ \gesperrt{Hat eine Abel'sche Gruppe $S$ einen Theiler $T$ vom
Grade $t$ und vom Index $j$, so giebt es unter den Charakteren von
$S$ genau $j$ und nicht mehr, die für alle Elemente von $T$ den
Werth 1 haben, während für jedes nicht in $T$ enthaltene Element
von $S$ wenigstens einer von ihnen von 1 verschieden ist. Diese
Charaktere bilden eine mit $S/T$ isomorphe Gruppe.}

Wir wollen diese Charaktere \gesperrt{zur Gruppe $T$ gehörig} nennen.

Da jeder Charakter $\chi_B$ einem bestimmten Elemente $B$ von $S$
entspricht, so wird auch, wenn $\chi_B$ die Gruppe der zu $T$
gehörigen Charaktere durchläuft, $B$ wegen der Formel \S.\,11, (18)
eine Gruppe durchlaufen, die vom Grade $j$ und mit der Gruppe der
$\chi_B$ und also auch mit der Gruppe $S/T$ isomorph ist. Diese
%CONT>>
